\section{Finite ordinals}

\import{ordinal.tex}
\import{order/finite.tex}

\begin{classification}\label{natural_number_is_ordinal}
    Every natural number is an ordinal.
\end{classification}
\begin{proof}
    Let $C=\{n\in\naturals\mid \text{$n$ is an ordinal}\}$.
    $\emptyset\in C$ by \obv{omega_empty,emptyset_is_ordinal}.
    For all $n\in C$ we have $\suc{n}\in C$
        by \obv{omega_successor,suc_ordinal}.
    Thus $C$ is an inductive set by \obv{inductive_set}.
    Follows by \auto{omega_minimal}.
\end{proof}

\begin{theorem}\label{natural_additive_difference}
    Let $m,n$ be natural numbers.
    Suppose $m\subseteq n$.
    Then there exists a natural number $k$ such that $n=m+k$.
\end{theorem}
\begin{proof}
    Let $D=\{v\in\naturals\mid
        \text{for every $u\in\naturals$ such that $u\subseteq v$
        there exists $k\in\naturals$ such that $v=u+k$}\}$.
    $\emptyset\in D$
        by \auto{omega_empty,subseteq_emptyset_iff,num_addition_zero_right}.
    Show that for every $v\in D$ we have $\suc{v}\in D$.
    \begin{subproof}
        Fix $v\in D$.
        $\suc{v}$ is a natural number by \obv{omega_successor}.
        Show that for every $u$ such that
            $u\in\naturals$ and $u\subseteq\suc{v}$
            there exists $k\in\naturals$ such that
            $\suc{v}=u+k$.
        \begin{subproof}
            Fix $u$.
            Assume $u\in\naturals$ and $u\subseteq\suc{v}$.
            \begin{byCase}
                \caseOf{$u=\suc{v}$.}
                    $\suc{v}=u+\emptyset$
                        by \obv{num_addition_zero_right}.
                    Follows by \auto{omega_empty}.
                \caseOf{$u\neq\suc{v}$.}
                    $u\in\suc{v}$
                        by \obv{ordinal_subseteq_implies_elem_or_eq,natural_number_is_ordinal,omega_successor}.
                    Hence $u=v$ or $u\in v$ by \obv{suc_iff}.
                    Hence $u\subseteq v$
                        by \obv{subseteq_refl,ordinal_elem_implies_subseteq,natural_number_is_ordinal}.
                    Take a natural number $k$ such that $v=u+k$.
                    $\suc{k}$ is a natural number by \obv{omega_successor}.
                    $\suc{v}=u+\suc{k}$
                        by \auto{num_addition_successor_right}.
                    Follows by \auto.
            \end{byCase}
        \end{subproof}
    \end{subproof}
    Hence $D$ is an inductive set by \obv{inductive_set}.
    Hence every natural number is an element of $D$
        by \obv{omega_minimal}.
    Thus $n\in D$.
\end{proof}

\begin{theorem}\label{natural_subseteq_iff_additive_difference}
    Let $m,n$ be natural numbers.
    Then $m\subseteq n$ iff there exists a natural number $k$ such that
        $n=m+k$.
\end{theorem}
\begin{proof}
    Show that if $m\subseteq n$, then there exists a natural number $k$
        such that $n=m+k$.
    \begin{subproof}
        Assume $m\subseteq n$.
        Follows by \obv{natural_additive_difference}.
    \end{subproof}
    Show that if there exists a natural number $k$ such that $n=m+k$,
        then $m\subseteq n$.
    \begin{subproof}
        Assume there exists a natural number $k$ such that $n=m+k$.
        Take a natural number $k$ such that $n=m+k$.
        $m\subseteq m+k$ by \obv{natural_addition_left_subseteq_sum}.
    \end{subproof}
    Follows by \auto.
\end{proof}

\subsection{Natural subtraction}

For natural numbers, subtraction is the unique additive difference when the
subtrahend is no larger than the minuend.  Taking the union of all such
differences makes the operation a set term without choosing a witness; the
laws below use it only in its ordinary range.

\begin{definition}\label{natural_difference}
    $\naturalDifference{n}{k}
        =\unions{\{d\in\naturals\mid n=k+d\}}$.
\end{definition}

\begin{theorem}\label{natural_difference_specification}
    Let $n,k$ be natural numbers.
    Suppose $k\subseteq n$.
    Then $\naturalDifference{n}{k}$ is a natural number and
        $n=k+\naturalDifference{n}{k}$.
\end{theorem}
\begin{proof}
    Take a natural number $d$ such that $n=k+d$ by
        \obv{natural_additive_difference}.
    Let $D=\{r\in\naturals\mid n=k+r\}$.
    Show that $D=\{d\}$.
    \begin{subproof}
        Show that every element of $D$ is equal to $d$.
        \begin{subproof}
            Fix $r\in D$.
            Hence $r$ is a natural number and $n=k+r$.
            Hence $k+r=k+d$.
            Follows by \obv{natural_addition_cancel_left}.
        \end{subproof}
        $d\in D$.
        Follows by \auto{setext,singleton_iff}.
    \end{subproof}
    $\unions{D}=d$ by
        \auto{setext,unions_iff,singleton_iff}.
    Hence $\naturalDifference{n}{k}=d$ by
        \obv{natural_difference}.
\end{proof}

\begin{proposition}\label{natural_difference_eq}
    Let $n,k,d$ be natural numbers.
    Suppose $k\subseteq n$ and $n=k+d$.
    Then $\naturalDifference{n}{k}=d$.
\end{proposition}
\begin{proof}
    $n=k+\naturalDifference{n}{k}$ and
        $\naturalDifference{n}{k}$ is a natural number by
        \obv{natural_difference_specification}.
    Follows by \obv{natural_addition_cancel_left}.
\end{proof}

\begin{proposition}\label{natural_difference_zero}
    Let $n$ be a natural number.
    Then $\naturalDifference{n}{\emptyset}=n$.
\end{proposition}
\begin{proof}
    $\emptyset\subseteq n$ by \obv{emptyset_subseteq}.
    $n=\emptyset+n$ by \obv{num_addition_zero_left}.
    Follows by \obv{natural_difference_eq,omega_empty}.
\end{proof}

\begin{proposition}\label{natural_difference_self}
    Let $n$ be a natural number.
    Then $\naturalDifference{n}{n}=\emptyset$.
\end{proposition}
\begin{proof}
    $n\subseteq n$ by \obv{subseteq_refl}.
    $n=n+\emptyset$ by \obv{num_addition_zero_right}.
    Follows by \obv{natural_difference_eq,omega_empty}.
\end{proof}

\begin{proposition}\label{natural_difference_successor_left}
    Let $n,k$ be natural numbers.
    Suppose $k\subseteq n$.
    Then
        $\naturalDifference{\suc{n}}{k}
            =\suc{\naturalDifference{n}{k}}$.
\end{proposition}
\begin{proof}
    Let $d=\naturalDifference{n}{k}$.
    $\naturalDifference{n}{k}$ is a natural number and
        $n=k+\naturalDifference{n}{k}$ by
        \obv{natural_difference_specification}.
    Hence $d$ is a natural number and $n=k+d$.
    Hence $\suc{n}=k+\suc{d}$ by
        \auto{num_addition_successor_right}.
    $n\subseteq\suc{n}$ by \obv{subseteq_self_suc_intro}.
    Hence $k\subseteq\suc{n}$ by \obv{subseteq_transitive}.
    $\suc{n}$ and $\suc{d}$ are natural numbers by
        \obv{omega_successor}.
    Follows by \obv{natural_difference_eq}.
\end{proof}

\begin{proposition}\label{natural_difference_successors}
    Let $n,k$ be natural numbers.
    Suppose $k\subseteq n$.
    Then
        $\naturalDifference{\suc{n}}{\suc{k}}
            =\naturalDifference{n}{k}$.
\end{proposition}
\begin{proof}
    Let $d=\naturalDifference{n}{k}$.
    $\naturalDifference{n}{k}$ is a natural number and
        $n=k+\naturalDifference{n}{k}$ by
        \obv{natural_difference_specification}.
    Hence $d$ is a natural number and $n=k+d$.
    Hence $\suc{n}=\suc{k}+d$ by
        \auto{num_addition_successor_left}.
    $k$ and $n$ are ordinals by
        \obv{natural_number_is_ordinal}.
    Hence $k\in n$ or $k=n$ by
        \obv{ordinal_subseteq_implies_elem_or_eq}.
    Hence $k\in\suc{n}$ by \auto{suc_iff}.
    $n\subseteq\suc{n}$ by \obv{subseteq_self_suc_intro}.
    Hence $k\subseteq\suc{n}$ by \obv{subseteq_transitive}.
    Hence $\suc{k}\subseteq\suc{n}$ by \obv{suc_subseteq_intro}.
    $\suc{n}$ and $\suc{k}$ are natural numbers by
        \obv{omega_successor}.
    Follows by \obv{natural_difference_eq}.
\end{proof}

\begin{proposition}\label{natural_nonzero_contains_zero}
    Let $n$ be a natural number.
    Suppose $n\neq\emptyset$.
    Then $\emptyset\in n$.
\end{proposition}
\begin{proof}
    $n$ is an ordinal by \obv{natural_number_is_ordinal}.
    Follows by \obv{ordinal_empty_or_emptyset_elem}.
\end{proof}

\begin{proposition}\label{natural_nonzero_successor_of_union}
    Let $n$ be a natural number.
    Suppose $n\neq\emptyset$.
    Then $\unions{n}$ is a natural number and
        $n=\suc{\unions{n}}$.
\end{proposition}
\begin{proof}
    Take a natural number $m$ such that $n=\suc{m}$
        by \obv{natural_zero_or_successor}.
    $m$ is \in-transitive
        by \obv{natural_number_is_ordinal,ordinal_is_transitiveset}.
    Hence $\unions{\suc{m}}=m$
        by \obv{transitiveset_iff_unions_suc}.
    Thus $\unions{n}=m$.
\end{proof}

\begin{proposition}\label{natural_successor_in_successor_bound}
    Let $i,r$ be natural numbers.
    Suppose $i\in r$.
    Then $\suc{i}\in\suc{r}$.
\end{proposition}
\begin{proof}
    $i$ is an ordinal and $r$ is an ordinal
        by \obv{natural_number_is_ordinal}.
    Hence $\suc{i}\subseteq r$ by \obv{ordinal_suc_subseteq}.
    $\suc{i}$ is an ordinal by \obv{suc_ordinal}.
    Hence $\suc{i}\in r$ or $\suc{i}=r$
        by \obv{ordinal_subseteq_implies_elem_or_eq}.
    Follows by \auto{suc_iff}.
\end{proof}

\begin{proposition}\label{naturals_is_transitive_set}
    $\naturals$ is \in-transitive.
\end{proposition}
\begin{proof}
    Let $C=\{n\in\naturals\mid n\subseteq\naturals\}$.
    $\emptyset\in C$ by \auto{omega_empty,emptyset,subseteq}.
    Show that for every $n\in C$ we have $\suc{n}\in C$.
    \begin{subproof}
        Fix $n\in C$.
        Then $n\in\naturals$ and $n\subseteq\naturals$.
        Hence $\suc{n}\in\naturals$ by \obv{omega_successor}.
        $\suc{n}\subseteq\naturals$ by \obv{suc_subseteq_intro}.
    \end{subproof}
    Hence $C$ is an inductive set by \obv{inductive_set}.
    Hence $\naturals\subseteq C$ by \obv{omega_minimal,subseteq}.
    Thus for every $n\in\naturals$ we have $n\subseteq\naturals$
        by \auto{subseteq}.
    Follows by \obv{transitiveset_iff_subseteq}.
\end{proof}

\begin{proposition}\label{omega_is_ordinal}
    $\naturals$ is an ordinal.
\end{proposition}
\begin{proof}
    Every element of $\naturals$ is an ordinal
        by \obv{natural_number_is_ordinal}.
    Follows by
        \obv{naturals_is_transitive_set,transitive_set_of_ordinals_is_ordinal}.
\end{proof}

\begin{proposition}\label{natural_number_is_finite}
    Every natural number is finite.
\end{proposition}
\begin{proof}
    Follows by \auto{finite,id_elem_bijections}.
\end{proof}

\begin{proposition}\label{omega_below_infinite_ordinal}
    Suppose $\alpha$ is an ordinal and $\alpha$ is infinite.
    Then $\naturals\subseteq\alpha$.
\end{proposition}
\begin{proof}
    Suppose not.
    $\naturals$ is an ordinal by \obv{omega_is_ordinal}.
    Hence $\alpha\subseteq\naturals$ by
        \obv{ordinal_subseteq_or_subseteq}.
    Hence $\alpha\in\naturals$ or $\alpha=\naturals$ by
        \obv{ordinal_subseteq_implies_elem_or_eq}.
    $\alpha\neq\naturals$.
    Hence $\alpha\in\naturals$.
    Hence $\alpha$ is finite by \obv{natural_number_is_finite}.
    Contradiction.
\end{proof}

\begin{proposition}\label{ordinal_membership_relation_transitive}
    Suppose $\alpha$ is an ordinal.
    Then $\memrel{\alpha}$ is transitive.
\end{proposition}
\begin{proof}
    Follows by \auto{transitive,memrel_elem_inspect,memrel_elem_intro,pair_eq_iff,elem_of_ordinal_is_ordinal,ordinal_transitivity}.
\end{proof}

\begin{proposition}\label{ordinal_membership_relation_connex}
    Suppose $\alpha$ is an ordinal.
    Then $\memrel{\alpha}$ is connex on $\alpha$.
\end{proposition}
\begin{proof}
    Follows by \obv{connex,memrel_elem_intro,elem_of_ordinal_is_ordinal,ordinal_elem_connex}.
\end{proof}

\begin{theorem}\label{finite_subset_of_naturals_bounded}
    Suppose $F\subseteq\naturals$ and $F$ is finite.
    Then there exists a natural number $n$ such that $F\subseteq n$.
\end{theorem}
\begin{proof}
    Let $C=\cons{\emptyset}{F}$.
    $\emptyset\in\naturals$ by \obv{omega_empty}.
    Hence $C\subseteq\naturals$ by \auto{cons_subseteq_intro}.
    $C$ is finite by \obv{finite_cons}.
    $\emptyset\in C$ by \auto{cons_iff}.
    Hence $C$ is inhabited.
    $\memrel{\naturals}$ is transitive
        by \obv{omega_is_ordinal,ordinal_membership_relation_transitive}.
    $\memrel{\naturals}$ is connex on $\naturals$
        by \obv{omega_is_ordinal,ordinal_membership_relation_connex}.
    Take $m\in C$ such that
        $m$ is a greatest element of $C$ with respect to
            $\memrel{\naturals}$
        by \obv{finite_greatest_element}.
    $m\in\naturals$ by \obv{elem_subseteq}.
    Hence $\suc{m}\in\naturals$ by \obv{omega_successor}.
    For every $k\in F$ we have $k\in\suc{m}$
        by
        \auto{cons_iff,greatest_element,memrel_elem_inspect,suc_iff}.
    Hence $F\subseteq\suc{m}$ by \obv{subseteq}.
\end{proof}

\begin{theorem}\label{finite_ordinal_is_natural_number}
    Suppose $\alpha$ is an ordinal and $\alpha$ is finite.
    Then $\alpha$ is a natural number.
\end{theorem}
\begin{proof}
    Suppose not.
    Let $B=\{\beta\in\suc{\alpha}\mid
        \text{$\beta$ is an ordinal and $\beta$ is finite and
            $\beta$ is not a natural number}\}$.
    $\alpha\in B$ by \obv{suc_intro_self}.
    Thus $B$ is inhabited.
    Take $\beta$ such that $\beta$ is an \in-minimal element of $B$
        by \obv{regularity}.
    $\beta$ is an ordinal and $\beta$ is finite and
        $\beta$ is not a natural number.
    $\beta\neq\emptyset$ by \obv{omega_empty}.
    Thus $\beta$ is inhabited by \obv{nonempty_inhabited}.

    $\memrel{\beta}$ is transitive
        by \obv{ordinal_membership_relation_transitive}.
    $\memrel{\beta}$ is connex on $\beta$
        by \obv{ordinal_membership_relation_connex}.
    Take $m\in\beta$ such that
        $m$ is a greatest element of $\beta$ with respect to $\memrel{\beta}$
        by \obv{finite_greatest_element,subseteq_refl}.

    $\beta\subseteq\suc{m}$
        by \auto{subseteq,greatest_element,memrel,suc_iff,pair_eq_iff}.
    $\suc{m}\subseteq\beta$
        by \obv{elem_of_ordinal_is_ordinal,ordinal_suc_subseteq}.
    Hence $\beta=\suc{m}$ by \obv{subseteq_antisymmetric}.

    $m$ is an ordinal by \obv{elem_of_ordinal_is_ordinal}.
    $m\subseteq\beta$ by \obv{ordinal_elem_implies_subseteq}.
    Thus $m$ is finite by \obv{subseteq_finite_is_finite}.
    $\beta\in\suc{\alpha}$.
    $\suc{\alpha}$ is an ordinal by \obv{suc_ordinal}.
    Hence $m\in\suc{\alpha}$ by \obv{ordinal_transitivity}.

    Show that $m$ is a natural number.
    \begin{subproof}
        Suppose not.
        Then $m\in B$.
        Contradiction by \auto{disjoint}.
    \end{subproof}
    Thus $\suc{m}$ is a natural number by \obv{omega_successor}.
    Contradiction.
\end{proof}

\begin{theorem}\label{natural_number_iff_finite_ordinal}
    $n$ is a natural number iff $n$ is finite and $n$ is an ordinal.
\end{theorem}
\begin{proof}
    Follows by \auto{natural_number_is_finite,natural_number_is_ordinal,finite_ordinal_is_natural_number}.
\end{proof}
